\DocumentMetadata{
  pdfversion=2.0,
  pdfstandard=ua-2,
  lang=en-US,
  tagging=on,
  tagging-setup={math/setup=mathml-SE,tabsorder=structure}
}
\documentclass[11pt,letterpaper]{article}
\usepackage[margin=0.68in,includefoot]{geometry}
\usepackage{newcomputermodern}
\usepackage{amsmath,amscd,mathtools}
\usepackage{tikz,adjustbox}
\providecommand{\square}{\mdlgwhtsquare}
\usepackage[hidelinks]{hyperref}
\hypersetup{
  pdftitle={Analysis First-Year Exam, January 2022, Part 1},
  pdfauthor={University of Florida Department of Mathematics},
  pdfsubject={Graduate mathematics examination},
  pdfdisplaydoctitle=true,
  bookmarksopen=true
}
\setcounter{secnumdepth}{0}
\setlength{\parindent}{0pt}
\setlength{\parskip}{0.28em}
\setlength{\emergencystretch}{3em}
\setlength{\leftmargini}{2.15em}
\setlength{\leftmarginii}{2.65em}
\setlength{\leftmarginiii}{3em}
\pagestyle{empty}
\raggedbottom
\allowdisplaybreaks
\NewTaggingSocketPlug{block/list/label}{examproblem}{%
  \tagstructbegin{tag=Lbl}\tagstructend%
  \tagstructbegin{tag=\LBody}%
  \tagstructbegin{tag=\ProblemHeadingTag,title-o={\ProblemHeadingText}}%
  \tagmcbegin{tag=\ProblemHeadingTag,actualtext={\ProblemHeadingText}}%
  #2%
  \tagmcend%
  \tagstructend%
}
\newenvironment{examproblems}[2]{%
  \def\ProblemHeadingTag{#2}%
  \AssignTaggingSocketPlug{block/list/label}{examproblem}%
  \begin{enumerate}%
  \setcounter{enumi}{#1}%
  \renewcommand{\labelenumi}{\textbf{\theenumi.}}%
  \setlength{\topsep}{0.35em}%
  \setlength{\partopsep}{0pt}%
  \setlength{\itemsep}{0.48em}%
  \setlength{\parsep}{0.18em}%
}{\end{enumerate}}
\newenvironment{examsubparts}{%
  \AssignTaggingSocketPlug{block/list/label}{default}%
  \begin{enumerate}%
  \setlength{\topsep}{0.18em}%
  \setlength{\partopsep}{0pt}%
  \setlength{\itemsep}{0.16em}%
  \setlength{\parsep}{0.1em}%
}{\end{enumerate}}
\newenvironment{examinstructions}{%
  \par\begingroup\itshape
}{%
  \par\endgroup\vspace{0.18em}
}
\makeatletter
\renewcommand\subsection{\@startsection{subsection}{2}{0pt}%
  {-1.25ex plus -.25ex minus -.1ex}%
  {0.55ex plus .1ex}%
  {\normalfont\large\bfseries}}
\renewcommand{\@maketitle}{%
  \newpage\null\vskip -1em%
  {\footnotesize\bfseries
    \href{https://www.ufl.edu/}{University of Florida}, \href{https://math.ufl.edu/}{Department of Mathematics}\par}%
  \vskip -0.5em%
  \tagstructbegin{tag=H1,title={Analysis First-Year Exam, January 2022, Part 1}}%
  \tagpdfparaOff%
  \tagmcbegin{tag=H1}%
    {\LARGE\bfseries\href{https://gma.math.ufl.edu/exams/analysis-first-year/}{Analysis First-Year Exam}, January 2022, Part 1\par}%
  \tagmcend%
  \tagpdfparaOn%
  \tagstructend%
  \vskip 0.25em%
  \tagmcbegin{artifact}%
    \hrule height 1pt%
  \tagmcend%
  \vskip 1em%
}
\makeatother
\title{Analysis First-Year Exam, January 2022, Part 1}
\author{}
\date{}
\begin{document}
\maketitle
\thispagestyle{empty}
\begin{examinstructions}
Answer FOUR questions in detail. State carefully any results used without proof.
\end{examinstructions}
\begin{examproblems}{0}{H2}
\def\ProblemHeadingText{Problem 1.}
\item
Let \((a_n)\) be a bounded real sequence. Let
\[
U=\{u:(\exists N)(\forall n\geq N)(a_n<u)\}
\]
 and
\[
L=\{\ell:(\exists N)(\forall n\geq N)(a_n>\ell)\}.
\]
 Prove that \(\sup L\) and \(\inf U\) exist and satisfy \(\sup L\leq\inf U\).
\def\ProblemHeadingText{Problem 2.}
\item
Let~\(X\) be the set of all bounded real-valued functions on the set~\(S\). Prove that the rule
\[
d(f,g)=\sup\{|f(s)-g(s)|:s\in S\}
\]
 defines a metric on~\(X\). Are closed and bounded subsets of~\(X\) compact when~\(S\) is finite? Explain.
\def\ProblemHeadingText{Problem 3.}
\item
This problem has two parts.
\begin{examsubparts}
\item[\textnormal{(i)}]
Prove that if \((x_n)_{n=0}^{\infty}\) converges to~\(x\) in a metric space then the set \(\{x_n:n\in\mathbb{N}\}\cup\{x\}\) is compact.
\item[\textnormal{(ii)}]
Prove that the function \(f:X\to Y\) between metric spaces is continuous if and only if its restriction \(f|_K\) to each compact subset~\(K\) of~\(X\) is continuous.
\end{examsubparts}
\def\ProblemHeadingText{Problem 4.}
\item
Let \(f:X\to Y\) be a map between metric spaces. Consider the statements: Prove or disprove each of the implications (i)~\(\Rightarrow\) (ii) and (ii)~\(\Rightarrow\) (i).
\begin{examsubparts}
\item[\textnormal{(i)}]
\(f\) maps Cauchy sequences to Cauchy sequences;
\item[\textnormal{(ii)}]
\(f\) is uniformly continuous.
\end{examsubparts}
\def\ProblemHeadingText{Problem 5.}
\item
The function \(f:(0,\infty)\to\mathbb{R}\) is differentiable and \(\lim_{x\to\infty}f'(x)=b\in\mathbb{R}\).
\begin{examsubparts}
\item[\textnormal{(i)}]
For \(h>0\) fixed, calculate the limit of \((f(x+h)-f(x))/x\) as \(x\to\infty\).
\item[\textnormal{(ii)}]
Determine~\(b\) if it is also known that \(f(x)\to a\in\mathbb{R}\) as \(x\to\infty\).
\end{examsubparts}
\end{examproblems}
\end{document}
